% =====================================================================
% conceptos.tex
% Responsabilidad única: presentar, en lenguaje sencillo, los dos
% conjuntos de datos y sus variables, el protocolo de partición y las
% métricas que se usan en los ejercicios 6 a 10. Los tamaños de
% partición son macros que escribe el notebook.
% =====================================================================
\subsection{Antes de empezar: datos, protocolo y métricas}

\noindent En esta parte pasamos de los datos de juguete a dos conjuntos
reales que vienen incluidos en scikit-learn, la biblioteca de Python con
la que hicimos todo el trabajo. El notebook que entregamos por separado se
ejecuta de principio a fin y produce todas las figuras, tablas y cifras
que citamos aquí, así que ningún número está copiado a mano.

\paragraph{Caso A: progresión de la diabetes (regresión).}
El conjunto \texttt{load\_diabetes} tiene $442$ pacientes con diabetes a
los que se les midieron diez variables al inicio de un estudio, y la
respuesta es una medida numérica de cuánto avanzó la enfermedad un año
después (entre más alto, peor). Como la respuesta es un número, este es un
problema de \emph{regresión}. Lo cargamos con \texttt{scaled=False} para
trabajar con las unidades originales y no con la versión ya transformada
que trae la biblioteca por defecto. La Tabla~\ref{tab:vars-diabetes}
explica cada variable.

\begin{table}[H]
\centering
\small
\begin{tabular}{l>{\raggedright\arraybackslash}p{11.5cm}}
\toprule
Variable & Qué mide \\
\midrule
\texttt{age} & Edad en años. \\
\texttt{sex} & Sexo, codificado como 1 o 2. \\
\texttt{bmi} & Índice de masa corporal: peso (kg) dividido por la estatura
(m) al cuadrado, en kg/m$^2$. \\
\texttt{bp} & Presión arterial promedio. \\
\texttt{s1} & Colesterol total en sangre. \\
\texttt{s2} & Colesterol LDL (el llamado \emph{colesterol malo}). \\
\texttt{s3} & Colesterol HDL (el llamado \emph{colesterol bueno}). \\
\texttt{s4} & Razón entre colesterol total y HDL. \\
\texttt{s5} & Logaritmo del nivel de triglicéridos en sangre. \\
\texttt{s6} & Nivel de glucosa (azúcar) en sangre. \\
\texttt{target} ($y$) & Progresión de la enfermedad un año después; la
documentación no da su unidad, así que hablamos de \emph{unidades de la
escala del target}. \\
\bottomrule
\end{tabular}
\caption{Variables del caso A (\texttt{load\_diabetes}).}
\label{tab:vars-diabetes}
\end{table}

\paragraph{Caso B: diagnóstico de cáncer de mama (clasificación).}
El conjunto \texttt{load\_breast\_cancer} tiene $569$ muestras de tejido
tomadas con una aguja fina de masas en el seno; cada muestra se fotografió
al microscopio y un programa midió la forma de los núcleos de las células.
La respuesta es si el tumor resultó maligno o benigno, así que es un
problema de \emph{clasificación}. La biblioteca codifica maligno como 0,
pero nosotros lo recodificamos como $1$ = maligno y $0$ = benigno, porque
por convención la clase $1$ (la \emph{positiva}) es la que nos interesa
detectar. Hay $30$ predictores que salen de medir diez características de
los núcleos (Tabla~\ref{tab:vars-cancer}) y resumir cada una de tres
maneras: el promedio sobre los núcleos de la imagen (prefijo
\texttt{mean}), su error estándar, que indica qué tanto varía entre núcleos
(sufijo \texttt{error}), y el promedio de los tres valores más grandes
(prefijo \texttt{worst}, \emph{el peor}).

\begin{table}[H]
\centering
\small
\begin{tabular}{l>{\raggedright\arraybackslash}p{11.5cm}}
\toprule
Característica & Qué mide del núcleo \\
\midrule
\texttt{radius} & Radio: distancia promedio del centro al borde. \\
\texttt{texture} & Textura: qué tan disparejos son los tonos de gris dentro
del núcleo. \\
\texttt{perimeter}, \texttt{area} & Perímetro y área del contorno. \\
\texttt{smoothness} & Suavidad: qué tanto varía el radio de un punto a otro
del borde. \\
\texttt{compactness} & Compacidad: $\text{perímetro}^2/\text{área}-1$; es
mínima para un círculo. \\
\texttt{concavity} & Concavidad: qué tan profundas son las partes hundidas
del contorno. \\
\texttt{concave points} & Puntos cóncavos: cuántas partes hundidas tiene el
contorno. \\
\texttt{symmetry} & Simetría del núcleo. \\
\texttt{fractal dimension} & Dimensión fractal: qué tan irregular o
\emph{dentado} es el borde. \\
\bottomrule
\end{tabular}
\caption{Las diez características del caso B; cada una aparece como
\texttt{mean}, \texttt{error} y \texttt{worst}, para $30$ predictores.}
\label{tab:vars-cancer}
\end{table}

\paragraph{Por qué partimos los datos en tres.}
Un modelo siempre se ve mejor sobre los datos con los que aprendió, igual
que un estudiante al que le ponen en el examen las mismas preguntas del
taller. Para saber cómo le irá con datos nuevos hay que guardar datos que
el modelo nunca haya visto. Por eso dividimos cada conjunto en tres partes
con papeles distintos:
\begin{itemize}
    \item \textbf{Entrenamiento} (60\,\%): con estos datos el modelo aprende
    sus parámetros (por ejemplo, los coeficientes $\beta$).
    \item \textbf{Validación} (20\,\%): con estos comparamos modelos y
    tomamos decisiones, como cuánto regularizar o qué umbral usar.
    \item \textbf{Prueba} (20\,\%): la abrimos una sola vez al final, en el
    Ejercicio~10, para medir de forma honesta el modelo ya elegido.
\end{itemize}
En la práctica apartamos primero el 20\,\% para prueba y luego el 25\,\% de
lo que quedaba para validación ($0.25\times0.8=0.2$), con
\texttt{random\_state=42} para que la partición sea siempre la misma. En el
caso B además \emph{estratificamos}, es decir, obligamos a que las tres
partes tengan la misma proporción de malignos que el conjunto completo.
Así quedaron $\NRegTr/\NRegVal/\NRegTest$ observaciones en el caso A y
$\NClfTr/\NClfVal/\NClfTest$ en el caso B (entrenamiento, validación y
prueba). Tampoco reentrenamos al final con entrenamiento más validación,
porque eso habría cambiado los modelos justo después de haberlos elegido.
La Figura~\ref{fig:protocolo} resume el protocolo.

\paragraph{Parámetros e hiperparámetros.}
Los \emph{parámetros} los aprende el modelo solo a partir de los datos (los
coeficientes de una regresión, las preguntas de un árbol). Los
\emph{hiperparámetros} los fijamos nosotros antes de ajustar y controlan
qué tan flexible puede ser el modelo: en la regresión logística, $C$
(cuánto se castigan los coeficientes grandes); en los árboles,
\texttt{max\_depth} (cuántas preguntas seguidas puede hacer) y
\texttt{min\_samples\_leaf} (cuántas observaciones debe tener como mínimo
cada grupo final). Los hiperparámetros los elegimos con validación.

\paragraph{Estandarizar y encadenar pasos.}
Varias técnicas, como la regularización, se ven afectadas por las unidades
de las variables (no es lo mismo medir en metros que en milímetros). Para
ponerlas a todas en la misma escala usamos \texttt{StandardScaler}, que
transforma cada variable en
\[
    z = \frac{x - \text{promedio de } x \text{ en entrenamiento}}
             {\text{desviación estándar de } x \text{ en entrenamiento}},
\]
de modo que $z$ dice a cuántas desviaciones estándar del promedio está cada
dato ($z=1$ significa una desviación por encima). El promedio y la
desviación se calculan solo con entrenamiento y luego se \emph{aplican} a
validación y prueba, para no filtrar información de los datos que deben
quedar ocultos. Esto lo garantiza un \texttt{Pipeline}, que encadena la
estandarización y el modelo como si fueran un solo paso.

\paragraph{Modelos de referencia.}
Para saber si un modelo aporta algo lo comparamos con uno que no mira los
predictores: \texttt{DummyRegressor} predice siempre el promedio de $y$ en
entrenamiento y \texttt{DummyClassifier} asigna a todos la misma
probabilidad (la proporción de malignos). Un modelo útil debe ganarle
claramente a esta referencia.

\paragraph{Cómo medimos el error.}
La Tabla~\ref{tab:metricas} reúne las métricas que usamos. En regresión,
$y_i$ es el valor real, $\hat y_i$ el predicho y $n$ el número de
observaciones de la partición que se evalúa. En clasificación, cada
predicción cae en una de las cuatro casillas de la \emph{matriz de
confusión}: verdadero positivo ($TP$, maligno detectado), falso positivo
($FP$, benigno marcado como maligno), falso negativo ($FN$, maligno que se
nos escapó) y verdadero negativo ($TN$, benigno bien identificado).

\begin{table}[H]
\centering
\small
\renewcommand{\arraystretch}{1.5}
\begin{tabular}{l>{\centering\arraybackslash}p{4.3cm}>{\raggedright\arraybackslash}p{7.2cm}}
\toprule
Métrica & Fórmula & Cómo se lee \\
\midrule
RMSE & $\sqrt{\frac1n\sum_i (y_i-\hat y_i)^2}$ & Error típico, en las
mismas unidades de $y$; castiga más los errores grandes. Menor es mejor. \\
MAE & $\frac1n\sum_i |y_i-\hat y_i|$ & Error promedio en valor absoluto,
en unidades de $y$. Menor es mejor. \\
$R^2$ & $1-\frac{\sum_i (y_i-\hat y_i)^2}{\sum_i (y_i-\bar y)^2}$ &
Fracción de la variabilidad de $y$ que explica el modelo; $0$ equivale a
predecir siempre el promedio y puede ser negativo si el modelo es peor que
eso. \\
\midrule
Precisión & $\frac{TP}{TP+FP}$ & De los que marcamos como malignos,
cuántos lo eran. \\
Recall (sensibilidad) & $\frac{TP}{TP+FN}$ & De los malignos reales,
cuántos detectamos. \\
F1 & $\frac{2\cdot\text{Precisión}\cdot\text{Recall}}{\text{Precisión}+\text{Recall}}$
& Resume precisión y recall en un número; solo es alto si ambos lo son. \\
ROC-AUC & Área bajo la curva ROC & Probabilidad de que, tomando un maligno y un benigno al
azar, el modelo le asigne mayor probabilidad al maligno. $0.5$ es azar y
$1$ es perfecto. No depende del umbral. \\
Brier & $\frac1n\sum_i (p_i-y_i)^2$ & Qué tan cerca están las
probabilidades $p_i$ de lo que pasó ($0$ o $1$). Menor es mejor. \\
Costo $L$ & $FP+5\,FN$ & Costo didáctico del enunciado: dejar escapar un
maligno pesa cinco veces más que una falsa alarma. \\
\bottomrule
\end{tabular}
\caption{Métricas de la Parte II.}
\label{tab:metricas}
\end{table}
